\section{Introduction}

Every major infrastructure investment creates byproducts. The US interstate highway system, designed primarily for automotive and freight movement, incidentally created a national network of rest areas, fueling facilities, and interchange nodes that later became the physical substrate for services their designers never anticipated: food distribution chains, emergency vehicle routing corridors, and telecommunications infrastructure using highway rights-of-way. The Interstate 2.0 upgrade program creates another class of byproduct: transit infrastructure.

Nine T1/T1 diamond hubs and approximately 50 T1/T2 regional stops, built primarily to harden freight network connectivity and eliminate single-point-of-failure interchange conditions, also create the physical conditions for intercity bus service to operate at levels of reliability currently available only in the Northeast Corridor. The access roads are built. The utility connections are established. The parking is funded by the freight program. The bus platform increment --- the covered waiting area, the ticketing kiosk, the passenger restroom upgrade --- costs \$1.8--2.5 million per hub.

\subsection{The US Intercity Bus Crisis}

US intercity bus service has contracted dramatically over the past four decades. The Motor Carrier Act of 1982 deregulated intercity bus service, triggering a wave of consolidation and route abandonment that left rural and small-metro communities without viable intercity bus connections \citep{Pisarski2006}. Greyhound Lines, which had operated as the de facto national intercity bus carrier since the 1920s, filed for bankruptcy in 2020 and subsequently eliminated service on approximately 1,200 route-segments, most of them in the South, Midwest, and rural West \citep{Amtrak2023}. FlixBus US, which entered the American market in 2018, has expanded aggressively on high-density corridors (NY--Philadelphia, LA--San Francisco, Dallas--Houston) but has not penetrated the rural and small-metro markets abandoned by Greyhound.

The result is an intercity mobility gap that disproportionately affects transit-dependent populations: households without personal vehicles, elderly residents who have stopped driving, low-income workers without access to air travel, and agricultural and service-sector workers in rural communities. The Census Bureau's American Community Survey (Table B08201) identifies approximately 9.7 million zero-vehicle households in non-urban counties \citep{ACS_B08201_2022}. These households have no viable intercity travel option in large parts of the country.

\subsection{The I2.0 Transit Opportunity}

The Interstate 2.0 diamond hub program creates the physical infrastructure that makes bus operator market entry viable on routes that were previously unserved because no adequate boarding facility existed. A bus operator cannot efficiently serve a general highway rest area: the geometry is not designed for coach maneuvering, the facilities are not designed for passenger amenity standards, and the parking is not designed for passenger egress. A T1/T1 diamond hub with a dedicated passenger platform is a different asset entirely.

This paper quantifies the transit opportunity created by I2.0:
\begin{enumerate}
  \item Nine T1/T1 diamond hubs and approximately 50 T1/T2 regional stops bring an estimated 12.4 million transit-dependent Americans within 30 miles of a potential intercity transit hub --- the geographic prerequisite for a connection they currently lack.
  \item The hub increment investment --- \$2 billion on a \$253 billion program --- serves transit-dependent travelers at \$161 per traveler served, compared to \$2,200 per traveler served for comparable standalone bus terminal construction: a 14-fold efficiency advantage.
  \item More than 40 origin-destination city pairs gain first-time intercity bus connectivity through the T1 hub network, concentrated in the Southeast, Gulf Coast, and Midwest corridors where rail does not serve and post-Greyhound bus service is absent.
  \item T1/T1 hub locations align disproportionately with communities scoring high on C3 (Economic Opportunity Access) in the ROUTE rubric, making the transit layer a significant equity multiplier for the freight investment.
\end{enumerate}

\subsection*{Contributions}
\begin{enumerate}
  \item We provide the first systematic cost-efficiency comparison between highway-embedded transit hubs and standalone urban bus terminal construction for a national-scale intercity bus program.
  \item We demonstrate that the rubric-driven site selection for freight infrastructure produces, as a byproduct, near-optimal siting for transit equity --- a result that follows from the shared determinants of freight intensity and economic opportunity access.
  \item We identify 40+ specific city pairs where the T1 hub network creates a viable bus operator market entry condition that does not currently exist, quantifying the demand, travel time, and current service gap for each pair.
  \item We recommend a mandatory transit design standard for all T1/T1 diamond hub construction, showing that the retrofit cost penalty for freight-only design is 3--5 times the incremental cost of transit-compatible design.
\end{enumerate}
